\chapter{Conclusiones y trabajo futuro}
\label{chap:conclusiones}

\section{Conclusiones}

En este trabajo se ha desarrollado un asistente para gestionar citas sanitarias
en un entorno simulado. Las tres arquitecturas permiten consultar, reservar,
reprogramar y cancelar citas, y configurar recordatorios sin enviarlos
realmente. Todas comparten los mismos controles de identidad, representación,
confirmación y escritura sobre una agenda basada en un subconjunto de FHIR R4.

La personalización usa las preferencias autorizadas para ordenar las opciones
disponibles y adaptar las respuestas. La persona puede corregir lo que ha
pedido o pedir que no se use su perfil. Ninguna preferencia da permisos ni crea
huecos nuevos. CitaLab permite conversar con cada variante y revisar las
llamadas, las decisiones de política y las operaciones realizadas.

La evaluación de 276 episodios con dos modelos no muestra una ventaja clara ni
estadísticamente significativa de la arquitectura multiagente. Con Groq, M5
obtiene la misma tasa de éxito que H y una algo mayor que S1 (43 frente a 42 de
46, aunque una de las unidades a su favor es un falso positivo del oráculo), y
con Gemini queda por debajo de S1 (43 frente a 46). En cambio, su coste es
mayor con los dos modelos: hace el doble de llamadas, usa de media unas 2,4
veces los tokens de H y tiene el mayor tiempo de inferencia
(Tablas~\ref{tab:resultados-config} y~\ref{tab:coste}, y
Figura~\ref{fig:coste-exito}). Al reunir los dos
modelos, H obtiene notas de calidad iguales o superiores a las de M5 en las
seis dimensiones, tanto en la revisión humana como en el jurado
(Tabla~\ref{tab:calidad}). Por tanto, en estas gestiones la complejidad añadida
no compensa su coste. La explicación más probable, a falta de una ablación, es
que los pasos decisivos (interpretar la petición y actualizar la tarea) son
comunes a M5 y H, y que los roles adicionales actúan después de ellos
(apartado~\ref{sec:discusion-coste}).

Las comprobaciones de seguridad del oráculo no detectaron ningún incumplimiento
en los 276 episodios: ninguna escritura se hizo sin confirmación, las citas de otras personas no cambiaron y no hubo
cambios en otras citas. En cambio, la comprobación de la redacción protege la
veracidad menos de lo que parece, ya que de los 32 borradores que sustituyó por
posibles anuncios solo tres anunciaban o prometían un efecto que no existía.
Además, que la confirmación sea correcta no garantiza que se haga lo que la
persona quería. En nueve episodios, una petición ambigua se convirtió en una
reserva nueva, propuesta con claridad, aceptada y autorizada, cuando la persona
quería mover su cita (apartado~\ref{sec:discusion-intencion}). Es decir, los
controles protegen la agenda, pero no garantizan que se entienda bien lo que se
pide.

La calidad depende de todos los pasos que van de la petición a la respuesta.
Hubo errores en la interpretación, en la actualización del estado, en los
valores por defecto de las herramientas, en las instrucciones que construye el
código, en la redacción y en los textos fijos
(apartado~\ref{sec:discusion-cadena}). Un ejemplo es la personalización con
Groq. El asistente ordenó bien los huecos según el perfil, pero como no explicó
el orden, el usuario simulado, que no conocía su preferencia, eligió otra
opción. Para personalizar no basta con ordenar bien, también hay que explicar
por qué se ofrece primero una opción.

El estudio de casos muestra además que varios fallos que parecían de
comprensión del modelo venían en realidad del código que lo rodea. S1 respondía
en español a quien escribía en inglés porque su código le indicaba ese idioma,
y con la corrección respondió en inglés en los doce primeros turnos probados con
cada modelo. El filtro arrastrado y la consulta que no encontraba citas se
corrigen cambiando pocas líneas, y una indicación en la redacción bastó para
que Groq mantuviera la explicación del orden en las doce pruebas. En cambio, la
confusión entre cambiar una cita y reservar otra no se resuelve solo con
instrucciones. Una indicación al agente de contexto la corrigió con Gemini,
pero con Groq solo en una de cinco repeticiones
(apartado~\ref{sec:discusion-correcciones}).

Combinar varios instrumentos de evaluación ha sido imprescindible, porque cada
uno detectó algo que los demás no veían. El oráculo encontró fallos en los que
el usuario simulado terminó satisfecho. El jurado distinguió esos fallos casi
sin error y detectó un falso éxito que el oráculo dio por bueno. Las trazas
permitieron encontrar la causa de cada fallo. El jurado de dos LLM coincide de
forma moderada o buena con la revisión humana y es algo más severo
(Tabla~\ref{tab:juez-concordancia}). La estimación con PPI++, que se había
previsto como principal, se descartó porque la muestra revisada no representa
toda la campaña.

\section{Grado de cumplimiento de los objetivos}

Los objetivos del capítulo~\ref{chap:introduccion} se han cumplido dentro del
alcance de laboratorio:

\begin{enumerate}
  \item Se han modelado el dominio de citas y las cinco gestiones, con sus
        estados, propuestas y confirmaciones
        (capítulo~\ref{chap:dominio-regulado}).
  \item Los requisitos del RGPD, del Reglamento de IA y de la Ley 41/2002 se
        han convertido en reglas del motor de políticas y en pruebas técnicas,
        que se superan en la versión evaluada
        (Apéndice~\ref{app:requisitos-normativos}). Esto no equivale a un
        cumplimiento jurídico de un despliegue.
  \item Se han definido las responsabilidades, los contratos y las rutas de
        escalado de los cinco componentes de M5
        (capítulo~\ref{chap:diseno}).
  \item El simulador utiliza datos sintéticos y un subconjunto de recursos de
        FHIR R4 (capítulo~\ref{chap:implementacion}).
  \item Las tres variantes comparten herramientas y controles y difieren solo
        en cómo toman las decisiones semánticas
        (capítulos~\ref{chap:diseno} y~\ref{chap:implementacion}).
  \item Se han medido el éxito integral de extremo a extremo, la seguridad
        funcional, el cumplimiento de los controles, el uso de preferencias, la
        calidad de la conversación, el tiempo de inferencia y el consumo en una
        campaña completa con dos modelos (capítulo~\ref{chap:resultados}).
  \item El código, los escenarios, los manifiestos de la campaña, las
        evaluaciones y los programas que generan las tablas y repiten las
        comprobaciones están publicados en
        \url{https://github.com/Andros5/citalab} para repetir el análisis
        (apartado~\ref{sec:metodo-ejecucion}).
\end{enumerate}

Las comparaciones no permiten afirmar que una arquitectura tenga más éxito que
otra, un resultado que ya se contemplaba como posible en el
capítulo~\ref{chap:metodologia-evaluacion}.

\section{Limitaciones}

El corpus contiene conversaciones sintéticas escritas a partir de gestiones
conocidas. Aunque es distinto del que se usó para el ajuste, comparte sus
categorías y capacidades. El usuario simulado es otro modelo de lenguaje, más
cooperativo que una persona real. Aceptó operaciones que no correspondían a su
objetivo y, como no conocía su preferencia horaria, su elección dependió de
cómo se presentaban las opciones. Las agendas solo ofrecen citas presenciales y
no reproducen la complejidad de un centro sanitario.

Las familias y sus contrastes comparten información, lo que limita el número
de observaciones independientes y la capacidad de detectar diferencias
pequeñas. La revisión humana se ha hecho sin revisores externos y cubre una
muestra de 60 episodios de la primera ejecución, ninguno de ellos fallido. Por todo
ello, no se establece una clasificación definitiva de las arquitecturas ni se
pueden sacar conclusiones sobre la satisfacción de pacientes, la seguridad
clínica o el cumplimiento jurídico de un despliegue real.

El oráculo tampoco comprueba de forma exhaustiva el alcance administrativo del
discurso. El caso N-B05 incluye recomendaciones clínicas pese a obtener un
resultado automático satisfactorio (apartado~\ref{sec:limite-alcance-clinico}).
Por tanto, superar las comprobaciones materiales no equivale a respetar todos
los límites de la respuesta.

\section{Trabajo futuro}

\subsection{Mejoras del asistente}

Los errores observados indican las primeras mejoras del asistente, varias de
ellas ya probadas sin repetir la campaña (apartado~\ref{sec:discusion-correcciones}):

\begin{itemize}
  \item Detectar en S1 el idioma de cada mensaje, como ya hacen M5 y H.
  \item Distinguir en el estado de la conversación entre \enquote{sin cambios}
        y \enquote{sin filtro}, y hacer explícitos los valores por defecto de
        las herramientas.
  \item Mantener la explicación del orden de las opciones, con una indicación
        en la redacción o como texto fijo, y adaptar los textos fijos al
        tratamiento y al idioma del perfil.
  \item Sustituir la comprobación léxica de anuncios por un contraste con las
        citas y la propuesta del estado, que cubra también las promesas en
        futuro.
  \item Preguntar antes de proponer cuando una petición pueda ser una reserva
        nueva o el cambio de una cita existente del mismo servicio, teniendo en
        cuenta lo que dice el mensaje y no solo el servicio.
\end{itemize}

\subsection{Mejoras de la evaluación}

Para separar los factores que esta campaña no aísla, haría falta un estudio de
ablación con variantes concretas: M5 sin el agente de perfiles, S1 con la
consulta del perfil obligatoria desde el código, H sin el paso de redacción y
M5 con un agente de contexto distinto del de H. Estas variantes permitirían
comprobar las hipótesis del apartado~\ref{sec:discusion-coste}.

También sería útil probar los propios instrumentos con
episodios en los que se introducen errores conocidos (una reserva duplicada,
una promesa en futuro, un filtro arrastrado), para medir qué detectan el
oráculo, el jurado y el simulador \citep{doddapaneni2024fbi}. Una muestra
humana estratificada por bloque y por resultado del oráculo permitiría usar
PPI++ con garantías \citep{fisch2024stratifiedppi}. El usuario simulado
debería recibir también las preferencias de la persona que representa, con la
indicación de no decirlas si no se las preguntan, para que su elección refleje
lo que quiere y no solo lo que le explica el asistente.

\subsection{Hacia una puesta en producción}

CitaLab se ha evaluado en un laboratorio, con datos sintéticos y un usuario
simulado. Antes de ponerlo en producción habría que comprobar que funciona con
personas reales y que cumple los requisitos de un servicio sanitario. Estos
serían los pasos, en orden:

\begin{enumerate}
  \item \textbf{Estudio con usuarios reales.} Probar el asistente con
        pacientes, acompañantes y personal administrativo, usando perfiles
        sintéticos y con consentimiento informado. Además del éxito de cada
        gestión, se mediría la satisfacción y la facilidad de uso, con especial
        atención a las personas mayores o con poca alfabetización digital
        \citep{vandervaart2017dhl}, y se seguirían pautas de accesibilidad
        cognitiva \citep{w3c2021coga}.
  \item \textbf{Revisión independiente.} Varios revisores ajenos al
        desarrollo, a ser posible con experiencia en la gestión de citas,
        puntuarían las conversaciones con la misma rúbrica. Así se mediría el
        acuerdo entre revisores y se validaría el jurado también en los
        episodios fallidos.
  \item \textbf{Integración con sistemas reales.} Conectar el asistente con la
        agenda de un servicio de salud mediante una API FHIR, con
        identificación electrónica de las personas y comprobación real de la
        representación. Esto obligaría a revisar los permisos, la autenticación
        y la concurrencia, y a cumplir los esquemas nacionales de seguridad e
        interoperabilidad \citep{es2022rd311ens,es2010rd4eni}. En esta fase se
        podrían añadir las citas telefónicas o por videoconsulta y el envío
        real de los recordatorios.
  \item \textbf{Protección de datos y seguridad.} Con datos reales haría falta
        una evaluación de impacto en protección de datos \citep{eu2016gdpr} y
        decidir dónde se ejecutan los modelos. GPT-OSS 20B es un modelo
        abierto y podría ejecutarse en servidores propios, de modo que los
        datos no saldrían de la organización. También habría que probar la
        resistencia del asistente a instrucciones maliciosas incluidas en los
        mensajes o en los datos que recibe \citep{debenedetti2024agentdojo}.
  \item \textbf{Piloto y seguimiento.} Empezar por un centro y por las
        gestiones de menos riesgo, como consultar citas o configurar
        recordatorios, y derivar al personal administrativo lo que el asistente
        no pueda resolver. Durante el piloto se conservarían las trazas y se
        aplicarían el oráculo y el jurado a conversaciones reales anonimizadas,
        para detectar errores nuevos y cambios de comportamiento cuando se
        actualice un modelo.
\end{enumerate}

Si estos pasos salen bien, el asistente podría ampliarse a otros canales, como
el teléfono o la mensajería, con entrada por voz y en más idiomas, y a otras
gestiones administrativas. Cada ampliación tendría que evaluarse de nuevo con el
mismo método antes de ofrecerse a los usuarios.
